\chapter{คู่มือการปฏิบัติการและแบบจำลองการเทียบมาตรฐาน}
\label{chap:ch04}

\section*{แผนบริหารการสอนประจำบทที่ 4}
\addcontentsline{toc}{section}{แผนบริหารการสอนประจำบทที่ 4}

\begin{planbox}{หัวข้อเนื้อหาประจำบท}
\begin{enumerate}[leftmargin=1.5em, itemsep=1pt]
    \item สถาปัตยกรรมการสลับหน้าจอ 5 หน้าจอ (5-Screen Carousel UI)
    \item ระบบการวิเคราะห์ 3 โหมด: TinyML, แบบจำลองพหุนาม และเส้นโค้งมาตรฐานสด
    \item ทฤษฎีการถดถอยเชิงเส้นตรง OLS การประเมินค่า $R^2$, LOD และ LOQ
    \item การวิเคราะห์สมบัติเชิงแสงและความหนาแน่นของของเหลวตามเกณฑ์ ICUMSA
    \item ขั้นตอนการตรวจวัดตัวอย่างดินและของเหลวในภาคสนาม และการบันทึกข้อมูล 4 ไฟล์
\end{enumerate}
\end{planbox}

\begin{planbox}{วัตถุประสงค์เชิงพฤติกรรม}
\begin{enumerate}[leftmargin=1.5em, itemsep=1pt]
    \item อธิบายขั้นตอนการสลับหน้าจอและการควบคุมคำสั่งผ่านจอยสติ๊กได้อย่างถูกต้อง
    \item คำนวณค่าสถิติการถดถอย $m$, $c$, $R^2$, LOD และ LOQ ได้อย่างถูกต้องตามเกณฑ์ IUPAC
    \item ปฏิบัติการตรวจวัดธาตุอาหารดินและสมบัติของเหลวได้อย่างถูกต้อง แม่นยำ
\end{enumerate}
\end{planbox}

\begin{planbox}{กิจกรรมการเรียนการสอน}
\begin{enumerate}[leftmargin=1.5em, itemsep=1pt]
    \item การสาธิตการใช้งานเครื่องสเปกโทรโฟโตมิเตอร์ครบทั้ง 5 หน้าจอ
    \item ปฏิบัติการสร้างเส้นโค้งมาตรฐาน 5 จุดความเข้มข้นสารละลาย NPK
    \item ปฏิบัติการสแกนสเปกตรัมของเหลวตัวอย่างและวิเคราะห์ดัชนีหักเหและความหนาแน่น
\end{enumerate}
\end{planbox}

\begin{planbox}{สื่อการเรียนการสอน}
\begin{enumerate}[leftmargin=1.5em, itemsep=1pt]
    \item เครื่องสเปกโทรโฟโตมิเตอร์ Wio Terminal และการ์ด microSD
    \item ชุดสารละลายมาตรฐาน N, P, K ความเข้มข้น 0, 20, 40, 80, 160 mg/kg
    \item สารละลายตัวอย่างของเหลวและน้ำกลั่นบริสุทธิ์สำหรับเทียบศูนย์
\end{enumerate}
\end{planbox}

\begin{planbox}{การวัดและประเมินผล}
\begin{enumerate}[leftmargin=1.5em, itemsep=1pt]
    \item การประเมินทักษะการเทียบศูนย์และการปฏิบัติการสแกนสเปกตรัม
    \item การตรวจสอบความถูกต้องของไฟล์ข้อมูล NPK.csv, CALIB\_LOG.csv และ LIQUID\_LOG.csv
    \item การตรวจความถูกต้องของการตอบคำถามทบทวนท้ายบทที่ 4
\end{enumerate}
\end{planbox}

\clearpage

\section{ระบบควบคุมหน้าจอและการนำทาง}
การใช้งานเครื่องสเปกโทรโฟโตมิเตอร์ตรวจวัดธาตุอาหารในดิน NPK Level Meter 1.02 ได้รับการออกแบบสถาปัตยกรรมซอฟต์แวร์ให้รองรับระบบเมนู 5 หน้าจอหลัก (5-Screen Carousel UI) ดังแสดงในตารางที่ \ref{tab:controls_mapping}

\begin{table}[H]
\centering
\caption{ฟังก์ชันการควบคุมการนำทางและคำสั่งของตัวเครื่อง}
\label{tab:controls_mapping}
\begin{tabularx}{\textwidth}{p{3.5cm}p{2.2cm}X}
\toprule
\textbf{ปุ่มกด / จอยสติ๊ก} & \textbf{ตำแหน่ง} & \textbf{หน้าที่การทำงาน} \\
\midrule
จอยสติ๊กโยกซ้าย (\texttt{WIO\_5S\_LEFT}) & ด้านหน้า & สลับไปยังหน้าจอก่อนหน้า (Page Loop 1 $\leftrightarrow$ 5) \\
จอยสติ๊กโยกขวา (\texttt{WIO\_5S\_RIGHT}) & ด้านหน้า & สลับไปยังหน้าจอถัดไป \\
จอยสติ๊กโยกลง (\texttt{WIO\_5S\_DOWN}) & ด้านหน้า & หน้าสเปกตรัม: สั่ง \textbf{Auto-Scan} กวาดแสง \\
 & & หน้าคาลิเบรต: สั่ง \textbf{บันทึกจุดสารมาตรฐาน (Step Calib)} \\
 & & หน้าของเหลว: สั่ง \textbf{วิเคราะห์สมบัติของเหลว (Analyze)} \\
จอยสติ๊กโยกขึ้น (\texttt{WIO\_5S\_UP}) & ด้านหน้า & หน้า NPK: สลับโหมดคำนวณ \textbf{[AI] $\to$ [PL] $\to$ [SC]} \\
จอยสติ๊กกดตรงกลาง (\texttt{WIO\_5S\_PRESS}) & ด้านหน้า & หน้าคาลิเบรต / ของเหลว: สั่ง \textbf{Zero Blank ($I_0$)} เทียบน้ำกลั่น \\
 & & หน้าอื่น: เปิดหรือปิดสลับกันของ \textbf{แสงสีขาว (White LED)} \\
ปุ่ม C (ซ้ายสุดด้านบน) & ด้านบน & หน้าคาลิเบรต: เลือกคาลิเบรต \textbf{ไนโตรเจน (N - 465 nm)} \\
 & & หน้า NPK: เปิดหรือปิดสลับกันของ \textbf{แสงสีน้ำเงิน (Blue LED)} \\
ปุ่ม B (กลางด้านบน) & ด้านบน & หน้าคาลิเบรต: เลือกคาลิเบรต \textbf{ฟอสฟอรัส (P - 525 nm)} \\
 & & หน้า NPK: เปิดหรือปิดสลับกันของ \textbf{แสงสีเขียว (Green LED)} \\
ปุ่ม A (ขวาสุดด้านบน) & ด้านบน & หน้าคาลิเบรต: เลือกคาลิเบรต \textbf{โพแทสเซียม (K - 625 nm)} \\
 & & หน้า NPK: เปิดหรือปิดสลับกันของ \textbf{แสงสีแดง (Red LED)} \\
\bottomrule
\end{tabularx}
\end{table}

\section{ระบบประมวลผลการคำนวณ 3 รูปแบบ}
ในหน้าจอตรวจวัด NPK Level Meter ผู้ใช้งานสามารถกดโยกจอยสติ๊กขึ้นเพื่อสลับโหมดการคำนวณได้ 3 รูปแบบอย่างอิสระ
\begin{enumerate}[leftmargin=1.5em, itemsep=3pt]
    \item \textbf{โหมดปัญญาประดิษฐ์ [AI] (TinyML Deep Neural Network)} ใช้โครงข่ายประสาทเทียม Multi-Layer Perceptron ($6 \to 12 \to 8 \to 3$) ประมวลผลบนฮาร์ดแวร์ FPU พร้อมระบบชดเชยความขุ่นของสารละลายดิน (Turbidity Matrix Compensation)
    \item \textbf{โหมดพหุนามสากล [PL] (Classical Literature Polynomial)} คำนวณความเข้มข้นจากสัดส่วนแสงสัมพัทธ์ $r_p, g_p, b_p$ ตามแบบจำลองพหุนามดั้งเดิม
    \item \textbf{โหมดเส้นโค้งมาตรฐานสด [SC] (In-Situ Calibrated Standard Curve)} คำนวณความเข้มข้นจากสมการเส้นตรงที่ผู้ใช้วัดเทียบสดในแปลงเกษตรด้วยสารละลายมาตรฐาน
    \begin{equation}
        C_{\text{sample}} = \max\left(0.0, \, \frac{A_{\text{sample}} - c}{m}\right)
    \end{equation}
\end{enumerate}

\section{ทฤษฎีการถดถอยเชิงเส้นและมาตรวิทยาเคมีวิเคราะห์}
ระบบเฟิร์มแวร์ของ Wio Terminal ได้รับการพัฒนาให้มีระบบวิเคราะห์การถดถอยกำลังสองน้อยที่สุด (Ordinary Least Squares หรือ OLS) บนชิป SAMD51 โดยตรง จากคู่ลำดับความเข้มข้นและค่าการดูดกลืนแสงของสารมาตรฐาน 5 จุด ($C_i, A_i$) สำหรับ $C \in \{0, 20, 40, 80, 160\}\text{ mg/kg}$

\subsection{ความชัน จุดตัดแกน และสัมประสิทธิ์การตัดสินใจ}
\begin{align}
    m &= \frac{\sum_{i=1}^N (C_i - \bar{C})(A_i - \bar{A})}{\sum_{i=1}^N (C_i - \bar{C})^2} \\
    c &= \bar{A} - m \bar{C} \\
    R^2 &= \frac{\left[\sum_{i=1}^N (C_i - \bar{C})(A_i - \bar{A})\right]^2}{\sum_{i=1}^N (C_i - \bar{C})^2 \sum_{i=1}^N (A_i - \bar{A})^2}
\end{align}

\subsection{ขีดจำกัดการตรวจวัดและขีดจำกัดการหาปริมาณ}
ความคลาดเคลื่อนมาตรฐานของส่วนตกค้าง ($s_{y/x}$) คำนวณจาก
\begin{equation}
    s_{y/x} = \sqrt{\frac{\sum_{i=1}^N \left(A_i - (m C_i + c)\right)^2}{N - 2}}
\end{equation}

จากนั้นประเมินขีดจำกัดการตรวจวัด (Limit of Detection หรือ LOD) และขีดจำกัดการหาปริมาณ (Limit of Quantification หรือ LOQ) ตามเกณฑ์สากลของ IUPAC
\begin{align}
    \text{LOD} &= \frac{3.3 \cdot s_{y/x}}{m} \\
    \text{LOQ} &= \frac{10.0 \cdot s_{y/x}}{m}
\end{align}

\section{การตรวจวัดสเปกตรัมการดูดกลืนแสงหลายช่วงคลื่น}
ในหน้าจอที่ 3 ผู้ใช้งานสามารถกดโยกจอยสติ๊กลง (\texttt{WIO\_5S\_DOWN}) เพื่อสั่ง Auto-Scan ระบบจะทำการดับไฟ 50 ms เพื่ออ่านค่ากระแสมืด ($I_{\text{dark}}$) แล้วกวาดแสง 5 ช่วงความยาวคลื่น (465, 500, 525, 590, 625 nm) พร้อมพล็อตกราฟเส้น $A(\lambda)$ แสดงจุดมาร์กเกอร์และตำแหน่งพีคการดูดกลืนแสง ($\lambda_{\max}, A_{\max}$) หากค่าการดูดกลืนแสงสูงเกิน 1.50 ระบบจะแสดงสัญญาณเตือนความเข้มข้นสูงเกินย่านเชิงเส้นทันที

\section{การวิเคราะห์สมบัติเชิงทัศนศาสตร์และความหนาแน่นของของเหลว}
ในหน้าจอที่ 5 ผู้ใช้งานสามารถตรวจวิเคราะห์สมบัติเชิงแสงและกายภาพของตัวอย่างของเหลว (Liquid Optics \& Metrology) ได้อย่างละเอียด ประกอบด้วย
\begin{enumerate}[leftmargin=1.5em, itemsep=3pt]
    \item \textbf{ดัชนีหักเหของของเหลว (Refractive Index, $n$)} คำนวณจากความสัมพันธ์เชิงทัศนศาสตร์ของสารละลายตามมาตรฐานสากล ICUMSA
    \begin{equation}
        n = 1.3330 + 0.00143 \cdot (\text{Brix}) + 0.0000045 \cdot (\text{Brix})^2
    \end{equation}
    โดยที่ $n_{\text{water}} = 1.3330$ ที่อุณหภูมิห้อง 20 องศาเซลเซียส
    \item \textbf{ความหนาแน่นของของเหลว (Density, $\rho$)} คำนวณตามแบบจำลอง Lorentz-Lorenz และ Gladstone-Dale ในหน่วย $\text{g/cm}^3$
    \begin{equation}
        \rho = 0.9982 + 0.00385 \cdot (\text{Brix}) + 0.000012 \cdot (\text{Brix})^2
    \end{equation}
    \item \textbf{ความเข้มข้นสารละลายและน้ำตาล (Brix Degree, $^\circ\text{Bx}$)} ประเมินจากการลดทอนความเข้มแสงและการดูดกลืนแสงเฉลี่ย ($A_{\text{mean}}$)
    \begin{equation}
        \text{Brix} = A_{\text{mean}} \times 26.5
    \end{equation}
    \item \textbf{ความส่องผ่านแสง (\%T) และการจำแนกความขุ่น (Clarity)}
    ระบบจะจำแนกสถานะความขุ่นออกเป็น 3 ระดับ
    \begin{itemize}[leftmargin=1.5em, itemsep=2pt]
        \item \textbf{CLEAR (ใส)} เมื่อค่าการดูดกลืนแสงเฉลี่ย $A < 0.12$
        \item \textbf{SLIGHT TURBID (ขุ่นเบา)} เมื่อ $0.12 \le A < 0.38$
        \item \textbf{TURBID (ขุ่นแขวนลอย)} เมื่อ $A \ge 0.38$
    \end{itemize}
\end{enumerate}

\section{ขั้นตอนการปฏิบัติการตรวจวัดในภาคสนาม}
\subsection{ขั้นตอนการตรวจวัดธาตุอาหารดิน}
\begin{enumerate}[leftmargin=1.5em, itemsep=2pt]
    \item ชั่งตัวอย่างดินแห้ง 5 กรัม ผสมน้ำยาสกัดธาตุอาหาร เขย่า 5 นาที และตั้งทิ้งไว้ให้ตกตะกอนใส
    \item ดูดสารละลายใส 2 มิลลิลิตร ใส่หลอดทดลองและเติมน้ำยาทดสอบเฉพาะของ N, P หรือ K
    \item เทียบศูนย์แบลงก์ด้วยน้ำกลั่นในหน้า Calibrate โดยกดจอยสติ๊กตรงกลาง (\texttt{WIO\_5S\_PRESS})
    \item ใส่หลอดคิวเวตต์บรรจุสารตัวอย่างลงในช่องวัด ปิดฝากันแสง และอ่านค่า NPK บนหน้าจอ
\end{enumerate}

\subsection{ขั้นตอนการตรวจวัดสมบัติของเหลว}
\begin{enumerate}[leftmargin=1.5em, itemsep=2pt]
    \item เลื่อนไปหน้าจอที่ 5 (\texttt{PAGE\_LIQUID})
    \item บรรจุน้ำกลั่นบริสุทธิ์ลงในคิวเวตต์ ใส่ในช่องวัด และกดจอยสติ๊กตรงกลางเพื่อตั้งค่าอ้างอิงน้ำกลั่น
    \item เปลี่ยนเป็นตัวอย่างของเหลวที่ต้องการตรวจวัด และกดจอยสติ๊กลง (\texttt{WIO\_5S\_DOWN})
    \item อ่านค่าดัชนีหักเห $n$, ความหนาแน่น $\rho$, Brix และระดับความขุ่นบนหน้าจอ
\end{enumerate}

\section{สถาปัตยกรรมการบันทึกข้อมูลบน microSD Card}
ระบบทำการบันทึกข้อมูลลง microSD Card แบบ Quadruple Logging 4 ไฟล์หลัก
\begin{itemize}[leftmargin=1.5em, itemsep=2pt]
    \item \texttt{NPK.csv} บันทึกค่าความเข้มข้น N, P, K ต่อเนื่องทุกวินาที
    \item \texttt{SPECTRUM.csv} บันทึกค่า Absorbance สเปกตรัม 5 ย่านคลื่นเมื่อสั่ง Auto-Scan
    \item \texttt{CALIB\_LOG.csv} บันทึกประวัติสมการเส้นตรงมาตรฐาน $R^2$, ความชัน $m$, จุดตัด $c$, LOD, LOQ
    \item \texttt{LIQUID\_LOG.csv} บันทึกผลการวิเคราะห์สมบัติเชิงแสงและความหนาแน่นของเหลว
\end{itemize}

\section*{บทสรุปประจำบทที่ 4}
\addcontentsline{toc}{section}{บทสรุปประจำบทที่ 4}
ระบบซอฟต์แวร์ 5 หน้าจอของ Wio Terminal ผสานการทำงานระหว่างการตรวจวัดธาตุอาหารดินแบบเรียลไทม์ การสแกนสเปกตรัมแสง การสร้างกราฟมาตรฐาน OLS และการวิเคราะห์สมบัติของเหลวได้อย่างสมบูรณ์แบบ การรองรับโมเดลการคำนวณ 3 ระบบช่วยเพิ่มความยืดหยุ่นในการประยุกต์ใช้งาน ขณะที่การบันทึกข้อมูล 4 รูปแบบช่วยให้สามารถตรวจสอบย้อนกลับข้อมูลทางวิทยาศาสตร์ได้อย่างครบถ้วน

\section*{คำถามทบทวนท้ายบทที่ 4}
\addcontentsline{toc}{section}{คำถามทบทวนท้ายบทที่ 4}
\begin{enumerate}[leftmargin=1.5em, itemsep=3pt]
    \item จงเปรียบเทียบข้อดีและข้อจำกัดของโหมดคำนวณ [AI] และโหมดเส้นโค้งมาตรฐานสด [SC] ในการตรวจวัดดินภาคสนาม
    \item ค่าสัมประสิทธิ์การตัดสินใจ ($R^2$) มีความสำคัญอย่างไรในการสร้างเส้นโค้งมาตรฐานตามเกณฑ์ ISO/IEC 17025
    \item จงคำนวณค่า LOD และ LOQ หากการถดถอยเชิงเส้นตรงของสารมาตรฐานฟอสฟอรัสให้ความชัน $m = 0.00765$ และความคลาดเคลื่อน $s_{y/x} = 0.0065$
    \item จงอธิบายหลักการจำแนกระดับความขุ่นของของเหลว (Clarity) โดยใช้ค่าการดูดกลืนแสงเฉลี่ย
    \item เหตุใดการจัดเก็บข้อมูลแบบ Quadruple Logging จึงมีความสำคัญต่องานวิจัยการเกษตรแม่นยำ
\end{enumerate}
